\documentclass[11pt,a4paper]{article}

\usepackage[top=2.5cm,bottom=2.5cm,left=2.5cm,right=2.5cm]{geometry}
\usepackage{amsmath,amssymb}
\usepackage{graphicx}
\usepackage{booktabs}
\usepackage[numbers,sort&compress]{natbib}
\usepackage{xcolor}
\usepackage{microtype}
\usepackage[colorlinks=true,linkcolor=blue!50!black,citecolor=blue!50!black,urlcolor=blue!50!black]{hyperref}

\newcommand{\M}{\mathbf{M}}
\newcommand{\St}{\mathbf{S}}
\newcommand{\R}{\mathbb{R}}
\newcommand{\C}{\mathbb{C}}
\newcommand{\vhat}{\hat{\mathbf{v}}}
\newcommand{\w}{\mathbf{w}}
\newcommand{\vb}{\mathbf{v}}

\title{\textbf{Supplementary Material: Operator-Specific Forgetting in Matrix-Valued Associative Memory}}
\author{Javier Escobar\\
  Independent Researcher}
\date{\today}

\begin{document}
\maketitle

\section{A. Reproducibility manifest}\label{sec:A}

Reference commit: \texttt{cec4ffe}. Tag: \texttt{lab-radioactivo-v1}. Repository: \url{https://github.com/01javierescobar/operator-specific-forgetting}. Archived DOI: \href{https://doi.org/10.5281/zenodo.22051928}{10.5281/zenodo.22051928}.

The repository contains implementations, execution runners, evaluation probes, audit logs, model checkpoints, frozen claims, and execution transcripts for experimental milestones O01--O05. The primary trained-model gate is O05. Verification tooling:
\begin{itemize}
  \item \texttt{paper/verify\_numbers.py}: Recomputes all quantitative values and verification gates directly from frozen JSON artifacts.
  \item \texttt{paper/make\_figures.py}: Regenerates all publication figures (PDF and PNG at 300~DPI) from the same frozen artifacts.
\end{itemize}

\section{B. Notation and memory construction}\label{sec:B}

$D$ denotes the key and value dimension ($D \in \{64, 128\}$). $n$ is the number of stored key--value pairs. $c=(n-1)/D$ is the associative crosstalk load excluding the target item. $\mathbf{w}_i \in \C^D$ are fixed codebook keys with unit-modulus entries ($|\mathbf{w}_{i,k}|=1$, squared Euclidean norm $\|\mathbf{w}_i\|_2^2 = D$). $\mathbf{v}_i \in \R^D$ are real-valued target vectors. The memory state $\M \in \C^{D \times D}$ stores associations via matrix superposition:
\begin{equation}
\M = \sum_{i=1}^n \mathbf{w}_i \mathbf{v}_i^\top .
\end{equation}
The matched-filter query read using key $\mathbf{w}_j$ is computed as $\hat{\mathbf{v}}_j = \mathbf{w}_j^H \M / D$. The pooled residual metric is defined as
\begin{equation}
L = \frac{\sum_i \|\mathbf{r}_i\|^2}{\sum_i \|\mathbf{v}_i\|^2},
\end{equation}
where $\mathbf{r}_i$ is the post-erase query readout residual.

\section{C. Formal operator definitions}\label{sec:C}

\subsection{C.1 Key drift reread erase}
Memory items are stored unperturbed. At erase time, local oscillator drift alters the erase key to $\tilde{\mathbf{w}}_j = \mathbf{w}_j e^{i\delta}$. The erase operator re-reads using $\tilde{\mathbf{w}}_j$ and subtracts the outer product:
\begin{equation}
\hat{\mathbf{v}}_j = \frac{\tilde{\mathbf{w}}_j^H \M}{D}, \qquad \M' = \M - \tilde{\mathbf{w}}_j \hat{\mathbf{v}}_j.
\end{equation}
Because both the projection estimate and the subtractive unbinding use $\tilde{\mathbf{w}}_j$, the operation represents an exact orthogonal projection onto the subspace orthogonal to $\tilde{\mathbf{w}}_j$.

\subsection{C.2 State drift reread erase}
We use state drift to denote a write-time phase offset attached to the target association, followed by a mismatch between the stored phase coordinate used for re-reading and the reference coordinate used for unbinding. Specifically, the target association is stored as $(e^{i\delta} \mathbf{w}_j) \mathbf{v}_j^\top$. At erase time, the erase operator re-reads using the item's stored key coordinate $\tilde{\mathbf{w}}_j = e^{i\delta}\mathbf{w}_j$ to recover $\hat{\mathbf{v}}_j = \tilde{\mathbf{w}}_j^H \M / D$, but subtracts using the current reference clock key $\mathbf{w}_j$:
\begin{equation}
\M' = \M - \mathbf{w}_j \hat{\mathbf{v}}_j.
\end{equation}

\subsection{C.3 Real-truncated readout channel}
When evaluating the real channel, real truncation is applied to the re-read estimate before subtraction: $\hat{\mathbf{v}}_j \leftarrow \operatorname{Re}(\hat{\mathbf{v}}_j)$. At query time, the readout residual is evaluated as $\mathbf{r}_j = \operatorname{Re}(\mathbf{w}_j^H \M' / D)$.

\subsection{C.4 Representational erase control}
The erased value is re-presented from an external target rather than re-read from the state: $\M' = \M - \mathbf{w}_j \mathbf{v}_j^\top$. This serves as a control condition separating operator-level self-estimation from state residual dynamics.

\subsection{C.5 Normalized delta control}
The corrective delta-rule update modifies state $\St \in \R^{D \times D}$ via
\begin{equation}
\St_t = \St_{t-1} + \beta\, \mathbf{k}_t \bigl(\mathbf{v}_t - \mathbf{k}_t^\top \St_{t-1}\bigr)^\top .
\end{equation}
In the normalized-key delta control arm, keys are normalized ($\|\mathbf{k}_t\|_2^2 = 1$). For $0 < \beta < 2$, this maintains the directional feedback multiplier $|1 - \beta| < 1$ strictly inside the stable unit disk.

\section{D. Closed-form law derivations}\label{sec:D}

\subsection{D.1 Distributed crosstalk scaling}
For independent random unit-modulus keys $\mathbf{w}_i, \mathbf{w}_k \in \C^D$, the inner products satisfy $\mathbb{E}[|\mathbf{w}_i^H \mathbf{w}_k|^2] = D$ for $i \ne k$. After normalization by $D$, each interfering association contributes expected energy $\mathbb{E}[|\mathbf{w}_i^H \mathbf{w}_k / D|^2] = 1/D$. Summing over $n-1$ interfering associations yields the crosstalk load $c=(n-1)/D$, consistent with distributed associative memory capacity \citep{plate1995holographic,kanerva2009hyperdimensional,frady2018theory}.

\subsection{D.2 Key drift under complex channel}
The operator $\mathbf{P}_{\tilde{\mathbf{w}}_j}^\perp = \mathbf{I} - \tilde{\mathbf{w}}_j \tilde{\mathbf{w}}_j^H / D$ is an exact rank-1 projection matrix. When querying with reference key $\mathbf{w}_j = \tilde{\mathbf{w}}_j e^{-i\delta}$, the target component in $\M'$ along direction $\tilde{\mathbf{w}}_j$ is identically zero. Therefore, $L_{\text{key, complex}}(\delta) = 0$ in the algebraic model.

\subsection{D.3 State drift under complex channel}
With target item stored as $(e^{i\delta}\mathbf{w}_j)\mathbf{v}_j^\top$ and unbinding performed with reference key $\mathbf{w}_j$, the re-read estimate is:
\begin{equation}
\hat{\mathbf{v}}_j = \frac{\tilde{\mathbf{w}}_j^H \M}{D} = e^{-i\delta}\mathbf{v}_j^\top + \frac{e^{-i\delta}}{D}\sum_{k \ne j} \mathbf{w}_j^H \mathbf{w}_k \mathbf{v}_k^\top .
\end{equation}
Subtracting $\mathbf{w}_j \hat{\mathbf{v}}_j$ yields the post-erase state:
\begin{equation}
\M' = (e^{i\delta}-1)\mathbf{w}_j \mathbf{v}_j^\top + \sum_{k \ne j} \mathbf{w}_k \mathbf{v}_k^\top - \mathbf{w}_j \frac{e^{-i\delta}}{D}\sum_{k \ne j} \mathbf{w}_j^H \mathbf{w}_k \mathbf{v}_k^\top .
\end{equation}
Reading with reference key $\mathbf{w}_j$ yields target residual $(e^{i\delta}-1)\mathbf{v}_j^\top$ and interference residual
\begin{equation}
\mathbf{r}_{\text{cross}} = \bigl(1-e^{-i\delta}\bigr)\frac{1}{D}\sum_{k \ne j}(\mathbf{w}_j^H \mathbf{w}_k)\mathbf{v}_k^\top .
\end{equation}
Taking expected squared norms over independent unit-modulus keys gives
\begin{equation}
\mathbb{E}[\|\mathbf{r}_j\|^2] = |e^{i\delta}-1|^2 \|\mathbf{v}_j\|^2 + |1-e^{-i\delta}|^2 \frac{n-1}{D}\|\mathbf{v}_j\|^2 = 2(1-\cos\delta)(1+c)\|\mathbf{v}_j\|^2.
\end{equation}
Because the factor $2(1-\cos\delta)(1+c)$ is invariant across target vectors, the ratio of the sum of expected residual energies to the sum of target energies preserves the identical pooled law:
\begin{equation}
\frac{\sum_j \mathbb{E}[\|\mathbf{r}_j\|^2]}{\sum_j \|\mathbf{v}_j\|^2} = 2(1-\cos\delta)(1+c).
\end{equation}

\subsection{D.4 Real-truncated channel derivations}
Applying real projection $\operatorname{Re}(\cdot)$ both at the re-read estimate $\hat{\mathbf{v}}_j \leftarrow \operatorname{Re}(\hat{\mathbf{v}}_j)$ and at the final query readout $\mathbf{r}_j = \operatorname{Re}(\mathbf{w}_j^H \M' / D)$ projects out the imaginary component of the phase differences.
\begin{itemize}
  \item \textbf{Key drift (Re channel)}: Under key drift, the target contribution to the real query residual is $\sin^2\delta \cdot \mathbf{v}_j^\top$, yielding target energy $\sin^4\delta \|\mathbf{v}_j\|^2$. The interfering crosstalk from $n-1$ unperturbed associations integrates over independent random phases to $\frac{c}{4}(1-\cos 2\delta)\|\mathbf{v}_j\|^2$. Cross-terms between target and interference vanish in expectation due to zero-mean key correlations ($\mathbb{E}[\mathbf{w}_j^H \mathbf{w}_k]=0$ for $k \ne j$). Summing gives:
  \begin{equation}
  L_{\text{key, Re}}(\delta) = \sin^4\delta + \frac{c}{4}\bigl(1-\cos 2\delta\bigr).
  \end{equation}
  \item \textbf{State drift (Re channel)}: The target association written with phase offset $e^{i\delta}$ yields an unbinding target residual $(1-\cos\delta)\mathbf{v}_j^\top$, with energy $(1-\cos\delta)^2 \|\mathbf{v}_j\|^2$. The unbinding subtracts the real part of the re-read crosstalk, leaving residual interference energy $c(1-\cos\delta)\|\mathbf{v}_j\|^2$. Independent key phases ensure vanishing cross-terms in expectation, yielding:
  \begin{equation}
  L_{\text{state, Re}}(\delta) = (1-\cos\delta)^2 + c\,(1-\cos\delta).
  \end{equation}
\end{itemize}

\section{E. Write stability analysis}\label{sec:E}

Along direction $\mathbf{k}_t$, projecting the state error yields the first-order scalar recursion
\begin{equation}
s_t = \bigl(1 - \beta\|\mathbf{k}_t\|^2\bigr) s_{t-1} + \beta\, q_t,
\end{equation}
where $q_t$ represents cross-talk from other stored patterns. The homogeneous dynamics are stable if and only if $|1-\beta\|\mathbf{k}_t\|^2| < 1$, which requires $0 < \beta\|\mathbf{k}_t\|^2 < 2$.

With $\beta=1$ and unnormalized embedding initialization ($\|\mathbf{k}_t\|^2 \approx 8$), the multiplier $|1-8| = 7$ induces exponential growth ($\|\St_t\| \sim 7^t$), causing numerical explosion within early epochs ($\|\St\| \approx 10^{11}$--$10^{12}$). Key normalization ensures $\|\mathbf{k}_t\|_2^2 = 1$; for $0 < \beta < 2$, this yields $|1-\beta| < 1$ and keeps the directional dynamics inside the unit disk.

\section{F. Experimental protocol (O05 milestone)}\label{sec:F}

\begin{itemize}
  \item \textbf{Dimensions}: Primary scale $D=128$ ($n=40$ pairs, predicted load $c=39/128 \approx 0.3047$, five seeds) compared against $D=64$ ($n=20$ pairs, predicted load $c=19/64 \approx 0.2969$, three seeds), achieving approximately matched load (2.7\% difference).
  \item \textbf{Vocabulary}: Discrete alphabet of 105 tokens (48 key tokens, 48 value tokens, special marker tokens). The token vocabulary is used strictly to construct controlled key--value associative sequences; the evaluation is not intended as a natural-language modeling benchmark.
  \item \textbf{Arms}: \texttt{wave\_complex}, \texttt{wave\_re}, \texttt{normalized\_delta}.
  \item \textbf{Events and Units}: 320 forgetting events per arm and seed (exceeding pre-registered minimum of 300). For all inferential statistical tests, the random seed was treated as the independent experimental unit.
  \item \textbf{Drift grid}: $\delta \in \{0.0, 0.1, 0.2, 0.3, 0.4, 0.5\}$.
  \item \textbf{Ridge probe}: Linear Ridge regression ($\alpha=0.1$, within-model split of 75\% train [$n=240$] / 25\% test [$n=80$]) mapping readout residuals $\mathbf{r}_j$ to continuous targets $\mathbf{v}_j \in \R^D$. Continuous predictions were decoded to the discrete value vocabulary by nearest cosine distance, with exact match recorded when the recovered token matched the true target token.
  \item \textbf{Selectivity}: $\text{Selectivity} = \mathrm{EM}_{\text{alive}} - \mathrm{EM}_{\text{erased}}$.
  \item \textbf{Equivalence test}: Two one-sided $t$-test (TOST at $\alpha=0.05$, two-sided $90\%$ CI) on retrieval selectivity between \texttt{wave\_complex} and \texttt{normalized\_delta} with pre-registered margin $\varepsilon=0.02$.
\end{itemize}

\section{G. Numerical results summary}\label{sec:G}

At $D=128$, all 15 seed--arm combinations reached $\mathrm{EM}=1.000$; at $D=64$, all 9 seed--arm combinations likewise trained to $\mathrm{EM}=1.000$. The four closed-form laws predicted empirical pooled residuals with maximum absolute deviation of $0.0049$ at $D=128$ and $0.0094$ at $D=64$, remaining below the $0.02$ verification threshold. Dimension-matched load comparison passed for all four cells with maximum cross-scale difference $0.0064$.

In the key/complex condition, all 30 seed $\times$ drift runs produced algebraically zero and numerically null residuals below the numerical guard threshold $\varepsilon_{\text{null}} = 10^{-4}$. TOST yielded mean difference $-0.0031$ with $90\%$ CI $[-0.0206,\,0.0144]$; because the lower bound crosses $-\varepsilon=-0.02$ by $0.0006$ under $n=5$ seeds, the result is reported as inconclusive with respect to formal equivalence.

\section{H. Metric audit}\label{sec:H}

A simple mean of per-example ratios $\frac{1}{N}\sum_i (\|\mathbf{r}_i\|^2 / \|\mathbf{v}_i\|^2)$ introduced severe estimation variance in trained models due to small target vector norms in the denominator. The primary metric was therefore standardized as the pooled ratio of sums $L = \sum_i \|\mathbf{r}_i\|^2 / \sum_i \|\mathbf{v}_i\|^2$, stabilizing the estimator.

When the post-erase residual is algebraically zero, computing erased-item exact match by taking an argmax over random noise is uninformative. In our evaluation protocol, residuals below $\varepsilon_{\text{null}} = 10^{-4}$ are recorded as numerically null residuals.

\section{I. Negative results and scope boundaries}\label{sec:I}

The vector-state geometry did not support the intended $D$-scaling law. The ideal wave erase and ideal delta erase were algebraically identical under normalized keys. The unnormalized $\beta=1$ corrective delta baseline diverged due to unnormalized key norms. Learned codebooks were not included in the closed-form theory. These results define the scope of the present analysis.

\bibliographystyle{plainnat}
\bibliography{references}

\end{document}